% Generated draft: requires mathematical and source review.
\documentclass[11pt]{article}
\usepackage[margin=25mm]{geometry}
\usepackage{fontspec}
\setmainfont{Latin Modern Roman}
\usepackage{amsmath,amssymb,mathtools,unicode-math}
\setmathfont{Latin Modern Math}
\usepackage{enumitem,needspace,xurl,hyperref}
\setlength{\parindent}{0pt}
\setlength{\parskip}{5pt}
\setlength{\emergencystretch}{3em}
\newcommand{\N}{\mathbb{N}}
\newcommand{\Z}{\mathbb{Z}}
\newcommand{\Q}{\mathbb{Q}}
\newcommand{\R}{\mathbb{R}}
\newcommand{\C}{\mathbb{C}}
\newcommand{\F}{\mathbb{F}}
\newcommand{\A}{\mathbb{A}}
\newcommand{\PP}{\mathbb{P}}
\newcommand{\E}{\mathbb{E}}
\newcommand{\eps}{\varepsilon}
\newcommand{\dd}{\,\mathrm{d}}
\begin{document}
\section*{UnsolvedMath ID: 800133}
\textbf{Source problem code:} catalogue entry
\subsection{Definitions and mathematical statement}
Let \(\zeta(s)\) denote the meromorphic continuation of the Riemann zeta function, initially defined by \(\zeta(s)=\sum_{n=1}^{\infty}n^{-s}\) for \(\operatorname{Re}(s)>1\). A nontrivial zero is a zero \(\rho\) with \(0<\operatorname{Re}(\rho)<1\). A positive ordinate is a number \(\gamma>0\) such that \(\zeta(\beta+i\gamma)=0\) for some \(\beta\in(0,1)\).

Prove that the Riemann zeta function has at least one irrational positive ordinate. Equivalently, prove
\[
\exists\,\beta\in(0,1)\;\exists\,\gamma\in\R\setminus\Q:\qquad \gamma>0\quad\text{and}\quad\zeta(\beta+i\gamma)=0.
\]
No Riemann hypothesis or linear independence hypothesis is imposed.

The source defines an ordinate of an \(L\)-function as a positive imaginary part of a nontrivial zero. As motivation, it states that LI implies that at most one positive ordinate can be rational, and expresses the expectation that even one rational positive ordinate is impossible for classical analytic \(L\)-functions. The supplied excerpt does not specify the precise formulation or scope of LI, or define the class of classical analytic \(L\)-functions. These contextual qualifications are not additional hypotheses of the stated problem.
\subsection{Short English statement}
Prove that some nontrivial zero of the Riemann zeta function has a positive imaginary part that is irrational. This asks for only one such ordinate and imposes no linear independence hypothesis.
\subsection{Sources}
\begin{itemize}
\item Local UnsolvedMath registry, permanent ID \(800133\), source code \texttt{catalogue entry}, title ``Comparative Prime Number Theory --- Problem 1.'' The supplied registry labels the problem open; this status is not independently verified.
\item Hamieh, Kadiri, Martin, and Ng, \emph{Comparative Prime Number Theory Problem List} (2024), Problem (\#1). Supplied source URL: \url{https://arxiv.org/abs/2407.03530}. The source attributes this problem to Greg Martin and Lior Silberman, University of British Columbia; the registry's proposer field lists Hamieh, Kadiri, Martin, and Ng.
\item Supplied provenance records access on 2026-10-16 and extraction by \texttt{arxiv\_latex}. This is supplied provenance, not a claim of independent source access. Upstream revision: \texttt{d371ee917ac9647fbb06cd6fc529deec6f3dda82}.
\end{itemize}
\end{document}
